\subsection{The local completion argument: setup and the first stage}
\label{sec:D-local-first-stage}

The next two statements reconstruct Sections 2--4 of the supplied
manuscript \citet{D_OpenAI2026}. They are a checked derivation of that
source, with the comparison maps and the treatment of every inverse
annulus made explicit. No priority is claimed. The conclusion of this
subsection is $L_\ell\subset H$; the passage from this field to the
whole separable closure requires the subsequent second stage.

\paragraph{Fixed fields and the valued point.}
Let $p$ be a prime, let $k/\Q_p$ be finite, and let $X/k$ be smooth,
proper, geometrically connected, and of genus at least two. Set
$K=k(X)$ and fix a separable closure $K^{\mathrm{sep}}$ containing an
algebraic closure $\bar k$ of $k$. Fix a real-valued nonarchimedean
field norm on $K$ extending the norm on $k$, normalized by
$|p|=p^{-1}$, and fix one prolongation to $K^{\mathrm{sep}}$.
Thus its support on $K$ is zero. All fields below have these specified
embeddings and norms. Put
\[
 K_0=\bar k(X),\qquad
 \widetilde K=\operatorname{compositum}_{V\to X}
                    k(V)\subset K^{\mathrm{sep}},
 \qquad E=\widehat{K^{\mathrm{sep}}},\qquad
 H=\widehat{\widetilde K}\subset E.
\]
In the compositum, $V$ runs through connected finite \'etale covers and all their
$K$-embeddings in the fixed separable closure. The displayed inclusion
of $H$ identifies it with the closure of $\widetilde K$ in $E$.
Finite constant-field covers show that $\bar k\subset\widetilde K$.
Consequently
\[
                     \Omega=\widehat{\bar k}\subset H
\]
is a specified complete algebraically closed subfield. Here we use
the usual fact that the completion of an algebraically closed
rank-one valued field is algebraically closed. For
$v=-\log_p|\cdot|$ one has $v(\Omega^\times)=\Q$: this holds on
$\bar k$, and a nonzero limit has the same norm as every sufficiently
close approximant, so completion adds no value.

The inclusion $K_0\hookrightarrow H$ defines an $H$-valued point of
$X_{\bar k}$ supported at its generic point. Together with the fixed
inclusion $\Omega\hookrightarrow H$, it defines a specified
$H$-valued point of $X_\Omega$. Denote its analytic image by
$\xi\in X_\Omega^{\rm an}$. We retain the full valued point,
including its isometric map of completed residue fields
$\mathcal H(\xi)\hookrightarrow H$. Its support on $X_\Omega$ can
be closed, even though its projection to $X_{\bar k}$ is generic;
in particular, $\xi$ can be an $\Omega$-rational point. It cannot
be any of the algebraic points in $X(\bar k)$.

For a finite extension $F/K_0$, normalization gives a finite map of
smooth proper curves over $\bar k$. Its \emph{horizontal indices}
mean the ramification indices at their closed points. These discrete
valuations are trivial on $\bar k$ and have residue characteristic
zero. Fix an odd prime $\ell\ne p$ and define, inside
$K^{\mathrm{sep}}$,
\[
 L_\ell=\operatorname{compositum}\{F: F/K_0\text{ finite Galois,
                       every horizontal index divides }\ell\}.
\]
Then $\widetilde K\subset L_\ell$. Indeed, after adjoining constants
the universal finite \'etale field is still $\widetilde K$: every
finite \'etale cover of $X_{\bar k}$ descends to a finite extension
of $k$, and composition with that constant-field cover gives a finite
\'etale cover of $X$.

\begin{claim}{D.15}{Proved; analytic patching, evaluation, and horizontal indices}
\label{clm:D15}
The following assertions hold in the preceding setup.
\begin{enumerate}[label=(\roman*)]
\item Let $f:C\to X_{\bar k}$ be finite and surjective, with $C$
smooth, proper, connected, of degree $d$ and branch set $B$.
Choose a rational local parameter $t_b$ at each $b\in B$.
There are disjoint open coordinate disks $B_b$ about these points,
avoiding $\xi$, such that the entire inverse disk is a disjoint
union of disks with equations $t_b=z_c^{e_c}$, where
$e_c=e(c/b)$. Suppose $g:Y\to X_\Omega$ is finite and surjective,
with $Y$ smooth, proper and connected, and that there are integers
$m_b$ divisible by all $e_c$ over $b$ and radii
$0<r_b<R_b$ in $|\Omega^\times|$ strictly inside $B_b$ for which
$g^*t_b$ has an analytic $m_b$th root on all of
$g^{-1}A_b$, where $A_b=\{r_b<|t_b|<R_b\}$.
Set
\[
 U=X_\Omega^{\rm an}\setminus
            \bigcup_{b\in B}\{|t_b|\le r_b\}.
\]
There is a connected finite \'etale surjective cover $V_\Omega\to Y$
and an analytic map
\[
     V_\Omega^{\rm an}\big|_U\longrightarrow C_\Omega^{\rm an}
                         \quad\hbox{over }X_\Omega^{\rm an}.
\]
The notation on the left is the inverse image of $U$ under the
composite map to $X_\Omega$.
\item Suppose that this composite $V_\Omega\to X_\Omega$ is
isomorphic over $X_\Omega$ to the base change of a finite map
$V\to X_{\bar k}$ from a smooth proper connected curve. If
$\bar k(V)$ has a $K_0$-embedding into $H$, then so does
$\bar k(C)$. The analytic exterior map need not descend to $\bar k$.
\item For every positive integer $n$, the condition that horizontal
indices divide $n$ is preserved by subextensions, finite composita,
Galois closures, normalized pullback along a finite map of smooth
proper curves, and composition with a finite \'etale upper cover.
In particular, a finite curve cover with horizontal indices dividing
$\ell$ has a $K_0$-embedding of its function field into $L_\ell$.
\end{enumerate}
\end{claim}
\begin{proof}
These are the assertions of Lemmas 3.1, 3.3, 3.4 and Proposition 3.2
of \citet{D_OpenAI2026}; we give the construction and the precise
uses of analytic comparison.

\paragraph{Local normalization, including indices divisible by $p$.}
At a point $c$ above $b$, choose an analytic parameter $u$ and write
$f^*t_b=u^{e_c}a(u)$, with $a(0)\ne0$. Choose an $e_c$th root of
$a(0)$ in $\Omega$. On a sufficiently small disk the inequality
\[
 v\bigl(a(u)/a(0)-1\bigr)>v_p(e_c)+\frac1{p-1}
\]
holds uniformly. The convergent expression
\[
 a(u)^{1/e_c}=a(0)^{1/e_c}
       \exp\bigl(e_c^{-1}\log(a(u)/a(0))\bigr)
\]
therefore supplies the desired unit root. The function
$z_c=u a(u)^{1/e_c}$ has nonzero derivative at $c$, so analytic
inversion gives a smaller disk on which it is a coordinate and
$f^*t_b=z_c^{e_c}$. No condition $p\nmid e_c$ was used.

Do this in disjoint disks at the finitely many points above $b$.
Their complement in the proper analytic curve $C_\Omega^{\rm an}$
is compact. Its image is closed and does not contain $b$.
A small enough base disk therefore has its \emph{entire} inverse
image in those chosen disks. Shrink the base disk once more so that
within each coordinate chart the inverse image is exactly the disk
specified by the bound on $|z_c|^{e_c}$. Shrinking disks form a
neighborhood basis at $b$. Hausdorffness, the finiteness of $B$, and
$\xi\notin X(\bar k)$ let us choose the $B_b$ disjoint from one
another and from $\xi$, with radii in $|\Omega^\times|$.

\paragraph{Patching every sheet.}
\label{clm:D15-patching}
The open sets
\[
 Y_U=g^{-1}U,
 \qquad Y_b=g^{-1}\{|t_b|<R_b\}
\]
cover $Y^{\rm an}$; different $Y_b$ are disjoint and
$Y_U\cap Y_b=g^{-1}A_b$. On $Y_U$ take the pullback of
$C_\Omega^{\rm an}$, a finite \'etale cover of degree $d$ because
all branch points have been removed. On each $Y_b$ take the trivial
cover of degree $d$. Write $q_b^{m_b}=g^*t_b$ on $g^{-1}A_b$.
For every point $c$ above $b$, the corresponding algebra on this
overlap splits as
\[
 Z^{e_c}-g^*t_b
   =\prod_{\zeta\in\mu_{e_c}}
              \bigl(Z-\zeta q_b^{m_b/e_c}\bigr).
\]
The function $q_b$ is a unit there. For distinct roots of unity,
$(\zeta-\zeta')q_b^{m_b/e_c}$ is an invertible analytic function.
This remains true when $p\mid e_c$: its constant factor is nonzero
in the field $\Omega$, independently of its absolute value.
The Chinese remainder theorem thus trivializes every local factor.
Since $\sum_{c\mid b}e_c=d$, it trivializes all $d$ sheets on the
\emph{whole} overlap.

Choose these trivializations and glue the finite locally free
analytic algebras. There are no overlaps between distinct $Y_b$,
so no further triple-overlap compatibility is needed. The resulting
coherent algebra is finite \'etale locally on $Y^{\rm an}$.
For the proper $\Omega$-scheme $Y$, the coherent-sheaf equivalence
of \citet[Appendix A, Theorem A.1(ii)]{D_Poineau2010} algebraizes
this sheaf. Full faithfulness algebraizes its unit and
multiplication and preserves the algebra identities. Its relative
spectrum is finite over $Y$, and it is \'etale because
analytification detects that property
\citep[Fact 5.1.5]{D_Temkin2010}. Choose a nonempty connected
component $V_\Omega$. Its image in connected $Y$ is open and
closed, or equivalently its finite \'etale degree is a positive
locally constant function; hence it is surjective. It is smooth
and proper. The original projection from the pullback over $Y_U$
gives the claimed analytic exterior map. If $B=\varnothing$,
the pullback is already finite \'etale and only this connected
component step is needed.

\paragraph{Evaluation at the specified $H$-point.}
\label{clm:D15-evaluation}
Let $j:\bar k(V)\hookrightarrow H$ be the embedding in (ii).
It gives a generic $H$-valued point of $V$ over the fixed generic
point of $X_{\bar k}$. Together with $\Omega\hookrightarrow H$
it gives an $H$-point of $V\times_{\bar k}\Omega$, and the chosen
isomorphism over $X_\Omega$ transports it to $V_\Omega$.
Its image in $X_\Omega(H)$ is exactly the point defined by
$K_0\hookrightarrow H$ and $\Omega\hookrightarrow H$.
Its analytic image is consequently above $\xi\in U$.

For a complete nonarchimedean extension $H/\Omega$, algebraic and
analytic $H$-points agree
\citep[Fact 5.1.3]{D_Temkin2010}. More explicitly the lifted point
is an analytic point $v$ equipped with an isometric embedding
$\mathcal H(v)\hookrightarrow H$. If $c'$ is its image under the
exterior analytic map, that map supplies
\[
             \mathcal H(c')\longrightarrow\mathcal H(v)
                                      \longrightarrow H.
\]
Thus it supplies an $H$-point of $C_\Omega$ over the same
\emph{valued} point of $X_\Omega$, not merely over the same
underlying topological point. Projecting to $C$, the support lies
over the generic point of $X_{\bar k}$. Since $C\to X_{\bar k}$
is finite, that support is generic in $C$. We obtain a field
embedding $\bar k(C)\hookrightarrow H$ whose restriction to
$K_0$ is the specified inclusion. This proves (ii), without any
descent assertion for the exterior analytic map or its local
trivializations.

\paragraph{Horizontal ramification under the allowed operations.}
\label{clm:D15-horizontal}
At a horizontal point the completed base field is
$A=\bar k((t))$. A finite extension has residue field $\bar k$
and is $\bar k((u))$ for a uniformizer $u$: successive subtraction
of representatives from the coefficient field $\bar k$ gives the
power-series expansion in its complete valuation ring.
Writing $t=u^e a(u)$ with $a(u)$ a unit, Hensel's lemma extracts
an $e$th root of the unit, because this is a residue-characteristic
zero valuation. A change of parameter therefore gives $t=z^e$.
Inside a fixed algebraic closure the extension is
$A(t^{1/e})$, which is cyclic since $\bar k$ contains all roots
of unity. If $e\mid n$, this field lies in $A(t^{1/n})$.

Every completed local factor of a subextension, finite compositum,
or normal closure consequently remains inside that same cyclic
extension. This proves the index assertion for those operations.
For a base change write its local equation as $t=s^r a(s)$,
$r\ge1$. Extracting an $e$th root of the unit shows that each
normalized local factor has ramification index
$e/\gcd(e,r)$, which divides $e$. Finally indices multiply in a
tower and an \'etale upper cover contributes index one. Taking
the Galois closure of a cover whose indices divide $\ell$ and
embedding it in $K^{\mathrm{sep}}$ gives the final assertion.
This proves (iii). It concerns horizontal valuations, and does not
assert that composing two ramified covers preserves the bound.
\end{proof}

\begin{claim}{D.16}{Proved; roots on every inverse collar and $L_\ell\subset H$}
\label{clm:D16}
Let $\ell\ne p$ be the odd prime fixed above.
\begin{enumerate}[label=(\roman*)]
\item For a smooth proper connected curve $C/\Omega$ of genus at
least two, an annular coordinate $S$ on a node annulus of its stable
skeleton acquires an analytic $\ell$th root on the \emph{entire}
inverse annulus in a connected finite \'etale cover of $C$.
One cover can serve a finite collection of node annuli.
\item Let $C/\bar k$ have the same properness, smoothness and genus
hypotheses, and let $b_i$ be finitely many distinct points, with
rational local parameters $t_i$. Inside any sufficiently small
pairwise disjoint coordinate disks $B_i\subset C_\Omega^{\rm an}$
about them, there are concentric collars
\[
 A'_i=\{r_i<|t_i|<R_i\},\qquad
             r_i,R_i\in|\Omega^\times|,\quad 0<r_i<R_i,
\]
and a connected finite \'etale cover $Y\to C_\Omega$ on which
$t_i$ has an analytic $\ell$th root on all of the inverse image
of $A'_i$, for every $i$.
\item With the specified subfields of $E$ in the setup,
\[
                              L_\ell\subset H.
\]
\end{enumerate}
\end{claim}
\begin{proof}
This reconstructs Lemmas 4.1--4.3 and Proposition 4.4 of
\citet{D_OpenAI2026}. We first give the elementary analytic input,
then construct the covers, and finally apply D.15.

\paragraph{Units on complete open disks and annuli.}
An analytic unit on an open disk has constant norm in
$|\Omega^\times|$. On an open annulus $A=\{r<|S|<R\}$ every unit
has a factorization
\[
 u=cS^a(1+h),\qquad c\in\Omega^\times,\quad a\in\Z,
       \quad h\in\mathcal O(A),\quad |h(x)|<1\quad(x\in A).
 \tag{D.16.1}
\]
Here $cS^a$ is a single term of its Laurent expansion, and the
factorization holds on the entire open annulus.
To verify this, restrict to a circle $|S|=\rho$ with
$\rho\in|\Omega^\times|$. Its Laurent Gauss norm is
multiplicative. The norm of a nonzero analytic function is attained
by one coefficient term and belongs to $|\Omega^\times|$.
Normalize $u$ and $u^{-1}$ to norm one with reciprocal constants.
Their reductions are inverse units in
$\widetilde\Omega[s,s^{-1}]$, whose units are monomials.
Exactly one Laurent term of $u$ therefore dominates at this radius.
The dominant exponent cannot change as the radius moves through
the annulus. Indeed, convergence on any compact subannulus leaves
only finitely many candidate dominant terms. A change forces a tie
of two such terms; their coefficient valuations are rational, so
the tie radius lies in $|\Omega^\times|$, contradicting the
preceding uniqueness. One term dominates on every compact
subannulus, and hence on their union, giving (D.16.1).
The Gauss inequality also gives the strict bound on $h$ away from
the skeleton. For a disk the corresponding reduction ring on
closed subdisks is $\widetilde\Omega[s]$, with only constant
units. The same argument shows constant norm; evaluating at a
rational point puts that norm in $|\Omega^\times|$.

Because $1/\ell\in\Z_p$, the binomial series for
$(1+h)^{1/\ell}$ converges on every compact subannulus: the
continuous norm of $h$ has a maximum strictly below one there.
These roots agree on overlaps because they are given by the same
series. Also $c$ has a root in $\Omega$. Thus if $S$ acquires an
$\ell$th root on an analytic inverse image, every unit $u$ on $A$
does too, including when $a<0$. More generally the unit already
has a root on $A$ whenever $\ell\mid a$.

\paragraph{The exact lifting input.}
A split stable model of a proper curve determines a finite metric
graph, its skeleton, and at each vertex the smooth normalization
$D_v$ of the residue component, with a marked point for each
incident edge end. A loop has two distinct marks. The tube of a
node is the annulus retracting to the corresponding open edge.
For $g(v)=g(D_v)$, stability is
\[
                         2g(v)-2+\operatorname{val}(v)>0.
 \tag{D.16.2}
\]
The residue field $\widetilde\Omega$ is algebraically closed of
characteristic $p$, and all edge lengths lie in $\Q$.
We use \citet[Theorem 7.7]{D_ABBR2015} for this fixed target
curve. Its input is a tame covering of metrized complexes:
a finite harmonic graph map, finite tame residue-curve maps with
ramification indices equal to the edge expansion factors, all
fibres of all marked points represented by edge ends, and zero
augmented graph ramification away from punctures
\citep[Definitions 2.19 and 2.21]{D_ABBR2015}.
We take no punctures, so the lift has empty branch locus and is
finite \'etale; the lifted vertex set and skeleton are exactly the
inverse images of the target ones
\citep[Definitions 4.25 and 4.31]{D_ABBR2015}.
The formulation of Theorem 7.7 includes loop edges.

\paragraph{A nonseparating stable edge.}
Let $e$ be the edge for the given annulus $A$, and first suppose
that $e$ is nonseparating. Choose a simple circuit through $e$.
At a vertex on this circuit let $P,Q$ be its two circuit marks
on $D_v$. Multiplication by $\ell$ on $\operatorname{Pic}^0(D_v)$
is surjective over the algebraically closed residue field.
Thus choose a degree-zero divisor $T$ for which
$P-Q+\ell T$ is principal, say the divisor of $f_v$.
The normalized cover $z^\ell=f_v$ is connected, because its
order at $P$ is congruent to $1$ modulo $\ell$, and it has degree
$\ell$. It is totally ramified at $P,Q$ and unramified elsewhere.
Each other marked fibre therefore has $\ell$ distinct points.
At every vertex outside the circuit take $\ell$ identity copies
of its residue curve.

Over each circuit edge take one edge with expansion factor
$\ell$, and length the target length divided by $\ell$.
Over every other edge take $\ell$ edges of factor one, matching
all marked fibres at both ends. This prescription includes chords
between circuit vertices and both ends of a loop. At each circuit
vertex the total degree in every target direction is $\ell$;
at an identity-copy vertex it is one. The residue formula
\[
       2g(D'_v)-2=\ell(2g(D_v)-2)+2(\ell-1)
 \tag{D.16.3}
\]
accounts for exactly the two ramified circuit directions.
There is consequently zero augmented graph ramification.
All residue ramification is tame, since $\ell\ne p$.
The resulting covering is connected: lift a path from any vertex
outside the circuit until its first meeting with the circuit;
its endpoint is on the connected lifted circuit. The source is
stable as well. At a circuit vertex its valence is
$2+\ell(\operatorname{val}(v)-2)$ and (D.16.3) gives
\[
 2g(v')-2+\operatorname{val}(v')
       =\ell\bigl(2g(v)-2+\operatorname{val}(v)\bigr)>0.
\]
Identity copies preserve (D.16.2).

Lift this tame covering by the stated theorem. We explain why its
inverse over $A$ consists of the \emph{whole} source node annuli,
including all their off-skeleton points. An off-skeleton open ball
maps into the complement of the target skeleton, by the equality
of inverse skeleta. Its connected image lies in one target
off-skeleton ball. Continuity sends its attachment point to that
target ball's attachment point, since the latter is the unique
point of the target skeleton in the ball's closure. Thus the
retractions commute on these balls and on the skeleton.
The inverse of the retraction tube of the open edge $e$ is exactly
the union of the corresponding source open-edge tubes. These are
the full node annuli above $A$.
On each one the Laurent exponent of the pullback of $S$ is
$\ell$ or $-\ell$, by the prescribed expansion factor and the
choice of orientation. Formula (D.16.1) supplies an analytic
$\ell$th root there. The annuli are disjoint, so these functions
give a root on the entire inverse image of $A$.

\paragraph{A separating stable edge, with every lifted edge treated.}
Suppose now that $e$ separates the graph. Each of the two sides,
retaining the boundary half-edge, contains either a circuit or a
positive-genus vertex. Otherwise it would be a tree with $N$
genus-zero vertices and total valence $2N-1$, including that
half-edge. Stability requires total valence at least $3N$, a
contradiction.

Construct a connected degree-$\ell$ unramified covering of each
side. If it contains a circuit, take the cyclic graph cover from
a surjection of its graph fundamental group onto $\Z/\ell\Z$,
with identity residue maps. If instead a positive-genus vertex is
used, choose a line bundle $L$ of exact order $\ell$ on its residue
curve. Such bundles exist because its Jacobian has nonzero
$\ell$-torsion. A trivialization $L^{\otimes\ell}\simeq\mathcal O$
makes $\bigoplus_{j=0}^{\ell-1}L^{\otimes j}$ a cyclic algebra.
Its spectrum is finite \'etale, since $\ell$ is invertible,
and connected: the nontrivial summands have no global sections
as nontrivial degree-zero line bundles, so the algebra has global
sections only the residue field and admits no nontrivial idempotent.
Use this connected cover at the chosen vertex and $\ell$ identity
copies elsewhere, matching all marked fibres. Every copy connects
to the distinguished vertex by lifting a path to its first
occurrence, so this covered side is connected.

Each side has $\ell$ preimages of its boundary mark. Match these
in pairs by $\ell$ factor-one edges. Both covered sides are
connected, so each of the new edges above $e$ is nonseparating:
removing it leaves another edge joining the two sides. The residue
maps are unramified and all edge factors are one. Harmonicity
and zero augmented ramification follow directly, or from the
unramified Riemann--Hurwitz formula. At a vertex of local degree
$d$ one has
\[
 2g(v')-2+\operatorname{val}(v')
                =d\bigl(2g(v)-2+\operatorname{val}(v)\bigr)>0,
\]
so the source triangulation is stable. Lift the covering; the
result is finite \'etale and every annulus over $A$ is now a
nonseparating stable-node annulus.

Apply the preceding circuit construction to each of these
$\ell$ annuli separately. Take a nonempty connected component
of the finite fibre product of the resulting finite \'etale
covers over this intermediate curve. Its projection to each
connected factor is finite \'etale and nonempty, hence
surjective. Every previously constructed root persists by
pullback on its entire prescribed inverse annulus. On each
intermediate factor-one annulus the original $S$ is an analytic
unit; (D.16.1) gives its root once that annulus's coordinate
has a root. This proves (i) for the original separating edge.
Taking a further connected component of a finite fibre product
treats any finite list of original node annuli simultaneously.

\paragraph{Moving arbitrary small collars into stable node annuli.}
All the algebraic data $C,b_i,t_i$ in (ii) descend to a finite
extension of $k$. After another finite extension choose a split
semistable model with the $b_i$ as disjoint smooth sections.
Around the $i$th section let $u_i$ be an integral local parameter.
Successive blowups of its reductions produce node annuli
\[
                  |\pi|^{j+1}<|u_i|<|\pi|^j,
 \tag{D.16.4}
\]
where $\pi$ is a uniformizer of that finite local field.
On a sufficiently small disk $t_i/u_i$ is an analytic unit,
with constant norm by (D.16.1)'s disk assertion. Thus (D.16.4)
is also a concentric annulus for $t_i$. Taking $j$ large puts
it inside the originally prescribed $B_i$. Carry out these
modifications for all the sections and call the resulting
proper semistable model $\mathcal C$. Forget its markings.
The model need not be stable; its newly introduced nodes
are precisely why the following input is required.

We use the model-domination consequence of resolution of
nonsingularities in
\citet[Theorem A and Proposition 2.4(iv)]{D_MochizukiTsujimura2024}.
For a hyperbolic curve over a finite extension of $\Q_p$, the set
of all primes contains $p$ and has the cardinality required there;
its residue field is algebraic over $\F_p$. Applied to the
prescribed compactified model $\mathcal C$, those results give,
after a finite base extension, a connected finite \'etale cover
whose compactified stable model maps to $\mathcal C$.
Pass to split reduction and a geometric connected component.
After base change to $\Omega$ write the curve map as
$f:Y_0\to C_\Omega$, and the compatible model map as
$\mathcal Y_0\to\mathcal C$. We require a model morphism,
not a finite model morphism. The source stable model has only
finitely many nodes.

Fix one target node annulus $A_i\subset B_i$ from (D.16.4).
Choose an irrational real number $a$ strictly between its
additive endpoint valuations. We claim that every point above
the whole level $\{x\in A_i:v(t_i(x))=a\}$ lies in a source
node annulus whose node maps to the target node.
The specialization maps commute with the model morphism.
A source point reducing to a component generic point is the
corresponding model vertex, a type-2 point whose value group is
$|\Omega^\times|$. It gives rational valuation to $f^*t_i$,
so cannot lie above this level. A source point reducing to a
smooth closed point lies in its open disk tube. If any such
point lies above the level, its reduction maps to the target
node; therefore its \emph{entire} disk tube maps into $A_i$.
On that disk $f^*t_i$ is a unit, hence has constant norm in
$|\Omega^\times|$, again excluding the irrational level.
The remaining possible reductions are source nodes mapping
to the target node, proving the claim.

Let $N_i$ be the union of all these source node annuli. It is a
finite union of open sets. The complement
$Y_0^{\rm an}\setminus N_i$ is compact, since $Y_0$ is proper.
Choose a closed concentric subannulus $A_i^0\subset A_i$
containing the chosen level in its interior. On the compact set
\[
 (Y_0^{\rm an}\setminus N_i)\cap f^{-1}(A_i^0)
\]
the continuous real function $v(f^*t_i)$ has compact image
which avoids $a$. Choose an interval about $a$ with rational
endpoints, contained in the valuation interval of $A_i^0$ and
disjoint from that compact image. It defines a collar
$A'_i\subset A_i$ whose \emph{entire} inverse image is contained
in $N_i$. This compact-complement step controls every sheet.

Every source annulus in $N_i$ maps into $A_i$ by specialization
compatibility, so $f^*t_i$ is a unit on that full annulus.
Apply (i) to all source annuli in all the $N_i$ at once.
Their coordinates acquire $\ell$th roots in one connected
finite \'etale cover $Y\to Y_0$. The unit factorization
(D.16.1) then supplies roots of $f^*t_i$ on their entire inverse
images, and hence on all of the inverse image of $A'_i$.
The composite $Y\to C_\Omega$ is finite \'etale. This proves (ii).

\paragraph{The first-stage field inclusion, with the specified copy.}
Let $F/K_0$ be one of the finite Galois extensions defining
$L_\ell$, and let $C\to X_{\bar k}$ be its smooth proper curve
cover. If it is unramified then $F\subset\widetilde K\subset H$.
Otherwise choose the disjoint branch disks of D.15, avoiding
$\xi$. Assertion (ii), applied to $X_{\bar k}$, supplies a
connected finite \'etale $Y\to X_\Omega$ and collars on whose
entire inverse images the branch parameters have $\ell$th roots.
D.15(i), with $m_b=\ell$, yields a connected finite \'etale
$V_\Omega\to Y$ and its analytic exterior map to $C_\Omega$.
The composite $V_\Omega\to X_\Omega$ is finite \'etale.

For a proper connected scheme over an algebraically closed field,
extension to an algebraically closed overfield induces an
equivalence of finite \'etale cover categories
\citep[Expos\'e X, Corollary 1.8]{E_SGA1}. Apply this to
$\bar k\subset\Omega$: the composite descends, up to
isomorphism over $X_\Omega$, to a connected finite \'etale
$V\to X_{\bar k}$. Its function field has a $K_0$-embedding
into $\widetilde K$, by the definition and constant-field
description of that universal \'etale field. D.15(ii) evaluates
the exterior analytic map at this point and gives a
$K_0$-embedding $j:F\hookrightarrow H\subset E$.

It remains to identify its image with the previously specified
subfield $F\subset K^{\mathrm{sep}}\subset E$. Normality is
essential for this last statement: every minimal polynomial
over $K_0$ of an element of $F$ splits in the specified copy
of $F$. Hence $j$ sends every element to an element of that
copy, and $j(F)\subset F$. Equality follows from the equal
finite degrees. Since $j(F)\subset H$, the specified $F$
itself lies in $H$. Taking the compositum over these fields
proves $L_\ell\subset H$ inside $E$, as asserted.
\end{proof}
